\subsection*{\href{https://sourceacademy.org/sicpy/1.1.1\#p3}{Numbers}}

Python supports three numeric types: integers (\texttt{int}), floats
(\texttt{float}), and complex numbers (\texttt{complex}).

A numeric literal never includes a leading sign: a \texttt{+} or
\texttt{-} written before a number, as in \texttt{-2} or
\texttt{-3.14}, is unary-operator application (see ``Operator
Precedence and Associativity'' above), not part of the literal itself.
This is why \texttt{\textbf{-2 ** 2}} is \texttt{\textbf{-(2 ** 2)}}
(i.e.\ $-4$) rather than \texttt{\textbf{(-2) ** 2}} (i.e.\ $4$): the
sign is applied only after \texttt{\textbf{**}} has already bound the
unsigned literal \texttt{\textbf{2}} to its exponent.

\subsubsection*{Integers}

Integers can be represented in decimal notation. Additional base
notations are supported, such as binary (\texttt{0b1010} or
\texttt{0B1010}), octal (\texttt{0o777} or \texttt{0O777}), and
hexadecimal (\texttt{0x1A3F} or \texttt{0X1A3F}). Underscores
(\verb|_|) may be used for readability (e.g., \verb|1_000_000|).
Examples include \texttt{42}, \texttt{0b1101} (and, with a unary
minus applied, \texttt{-0b1101}), and \verb|0x_FF_00|.

\subsubsection*{Floats}
Floats use decimal notation with an optional decimal dot. Scientific
notation (multiplying the number by $10^x$) is indicated with the
letter \texttt{e} or \texttt{E}, followed by the exponent \texttt{x}.
Special values inf (infinity), -inf, and nan (not a number) are allowed.
Examples include \texttt{3.14}, \texttt{0.001e+05} (and, with a unary
minus applied, \texttt{-0.001e+05}), and \texttt{6.022E23}.

\subsubsection*{Complex Numbers}
A complex literal consists of a real part and an imaginary part,
written \texttt{<real>±<imag>j}, where \texttt{j} (or \texttt{J})
denotes the imaginary unit. Both parts are stored as floats. The
imaginary part is mandatory (e.g., \texttt{5j}, \texttt{0j}, and
\texttt{0+3j} is valid, \texttt{5} alone is real).
As with the sign of a real number, the \texttt{±} between the real and
imaginary parts is not part of a single token: only \texttt{<imag>j}
(e.g., \texttt{3j}) is a numeric literal, so \texttt{2+3j} parses as
the real literal \texttt{2} combined with the imaginary literal
\texttt{3j} by binary \texttt{+}, and \texttt{-4.5j} as unary
\texttt{-} applied to the imaginary literal \texttt{4.5j}.

Examples include \texttt{2+3j}, \texttt{-4.5J}, \texttt{0j}, and
\texttt{1e3-6.2E2J}.
